% ---------------------------------------------------------------------------
% Drop-in replacement for the strong- and weak-scaling subsections of the
% methods paper (Sec. 4), from the 2026-10 campaign on the JUPITER Booster
% (JSC): NVIDIA GH200, 4 GPUs/node, fp32, Pallas backend with the library
% default block shape (128 cells along the contiguous z axis, (1,1,128) on
% these grids), 1D x decomposition, one process per node driving its four
% GPUs (pytests/runners/_env.sh: sharded_layout), pairwise halo exchange
% (BackendConfig.halo_exchange = "pair_all_to_all").
% FD = WENO-Z (admissible face state, LSRK4; MHD with the staged Pallas CT
% kernels on several GPUs); FV = AthenaPK's VL2 (PLM, HLLC / HLLD + GLM).
% Every rung executes the compiled program twice on the same allocation;
% quoted runtimes are the warm second execution.  Multi-node rungs were
% serialized (one job on the network at a time, --dependency=singleton).
% Sharded runs are bit-identical to single-GPU runs of the same global grid
% (G = 2, 4 on one node, G = 8 on two nodes, all four configurations:
% data/weak_scaling/sharded_validation_pair_all_to_all_G00{2,4,8}.json).
%
% Data: examples/scripts/scaling/data/{weak_scaling,strong_scaling}/
%   tags gh200zpair (FD hydro), gh200zctpair (FD MHD, Pallas CT),
%   gh200znoctpair (FD MHD, default CT, single GPU), gh200vl2pair (FV)
% Regenerate:
%   python examples/scripts/scaling/make_weak_tables.py --figure --copy-to-paper
%   python examples/scripts/scaling/strong_scaling_speedup.py --plot --copy-to-paper
% The self-gravity table at the end is the earlier WENO-JS campaign
% (gh200fftw, one process per GPU) and has not been re-measured.
% ---------------------------------------------------------------------------

\subsection{Strong scaling}
\label{sec:strong_scaling}

\begin{figure}
    \centering
    \includesvg[width=\linewidth]{figures/strong_speedup_hydro_mhd.svg}
    \caption{Strong-scaling speedup $T(1)/T(G)$ versus resolution $N$
    (grid $2N\times N\times N$) on NVIDIA GH200 GPUs (4 per node), single
    precision, for the FD (WENO-Z) and FV (VL2) schemes on the Pallas
    backend; left hydrodynamics (sound wave), right MHD (CP Alfv\'en
    wave).  Line styles distinguish 4 GPUs (one node), 8 GPUs (two nodes)
    and 16 GPUs (four nodes).  Open markers: the single-GPU baseline does
    not fit into one GPU's memory and is extrapolated with a power law
    fitted to the three largest measured baselines.}
    \label{fig:strong_speedup}
\end{figure}

For the strong-scaling test the global grid $2N\times N\times N$ is fixed
and decomposed along $x$ over $G$ GPUs (Fig.~\ref{fig:strong_speedup}).
The setup, the measurement and the precision are those of the weak-scaling
test below (warm second execution, compilation excluded, one process per
node); the number of time steps grows as $N^{-3}$ below $N=512$ so that
every point performs the same work.  The single-GPU baselines are measured
up to $N=768$ for hydrodynamics and up to $N=512$ for MHD;
larger baselines are extrapolated (open markers).  For FD MHD the baseline
is the faster of the two single-GPU constrained-transport variants; the
multi-GPU runs use the Pallas constrained-transport kernels, which need
fewer halo exchanges.

Within one node (4 GPUs driven by one process) the speedup at the largest
measured baselines is $3.1$ (FD hydro, $N=768$), $3.6$ (FD MHD,
$N=512$), $2.8$ (FV hydro, $N=768$) and $3.1$ (FV MHD, $N=512$).  Across
nodes the speedup keeps growing with the GPU count: on 8 GPUs (two
nodes) FD reaches $6.2$ (hydro) and $6.8$ (MHD), FV $5.4$ and
$5.6$; on 16 GPUs (four nodes) FD reaches $11.6$ (hydro) and
$12.1$ (MHD), FV $9.7$ and $9.6$ (hydro at $N=768$, MHD at
$N=512$).  With extrapolated baselines at $N=1024$ the 16-GPU speedups
are $12.3$ (FD hydro), $14.0$ (FD MHD), $10.3$ (FV hydro) and
$12.1$ (FV MHD).  Small grids gain less, because their steps are not
much longer than the halo exchange: at $N=128$ the 8-GPU speedup is $3.4$
(FD hydro) and $2.3$ (FV hydro).

\subsection{Weak scaling}
\label{sec:weak_scaling}

For the weak-scaling tests the per-GPU sub-domain is held fixed while the
device count grows, so ideal scaling is a flat wall-clock curve.  The
initial condition is built in a globally sharded manner (each process
materializes only its local shard; for FD MHD this includes the staggered
vector-potential construction that keeps the discrete face-centered field
divergence-free), so the full grid never resides on a single host.  All
runs use single precision, the Pallas backend with blocks of 128 cells
along the contiguous $z$ axis, and a 1D $x$ decomposition on the JUPITER
Booster (NVIDIA GH200, 4 GPUs per node).  The FD scheme uses WENO-Z and
LSRK4; the FV scheme is the VL2 scheme with HLLC (hydro) or HLLD and GLM
cleaning (MHD).  Both schemes run on the same per-GPU sub-domain --
$128\times2048\times2048$ for hydrodynamics and $128\times2048\times1024$
for MHD -- so the tables compare them directly.  FD runs 10 fixed time
steps and FV, whose steps are an order of magnitude cheaper, 100, so that
both measure a similar wall-clock interval.  We quote per-GPU throughput in
Mcell/s/GPU: millions of cell updates (grid cells $\times$ time steps) per
second and per GPU.  The decomposed runs are bit-identical to single-GPU
runs of the same global grid (checked on 2, 4 and 8 GPUs for all four
configurations).

Each node runs one process that drives its four GPUs, so the halo exchanges
between the GPUs of a node stay within one process; with one process per
GPU the multi-node rungs are about $30\%$ slower (8 GPUs: $21.1$ versus
$14.4$~s for FD MHD, $16.5$ versus $11.8$~s for FD hydro, $19.5$ versus
$12.7$~s for FV hydro).  Compilation is excluded from all timings (the
program is compiled ahead-of-time before the timer starts), and one-time
first-launch costs are deduced explicitly: writing the wall-clock time of
$n$ time steps as $T(n)=W+n\,s$, with $W$ the one-time initialization
(chiefly NCCL communicator construction) and $s$ the steady per-step cost,
each rung executes the identical compiled program twice back-to-back on
the same node allocation.  The first (cold) execution measures $W+n\,s$,
the second (warm) execution $n\,s$; all quoted runtimes are the warm
execution.  The deduced $W$ is at most $1.7$~s at any rung.  All
multi-node rungs were run one at a time, since concurrent rungs on the
shared network inflate the inter-node cost.

The halos are exchanged in two phases over disjoint pairs of neighbouring
devices (even--odd, then odd--even including the periodic wraparound pair),
each an all-to-all within its pair, rather than by two collective permutes
over the whole ring of devices: the pairwise collectives run on two-device
communicators whose cost does not grow with the device count.  The data
moved is the same and the results are bit-identical; measured back to back
on the same allocation, the pairwise exchange is $5\%$ faster at 8 and 32
GPUs and $5$--$13\%$ faster at 128 GPUs (FD MHD $15.4\to13.4$~s, FV hydro
$15.6\to13.6$~s).  The per-pair communicators need device memory outside the
XLA memory pool, so all sharded runs reserve $10\%$ of the GPU memory for
NCCL.

\begin{table}
    \centering
    \small
    \setlength{\tabcolsep}{4pt}
    \caption{Weak scaling of the 3D sound-wave (hydro) test on NVIDIA GH200
    GPUs, FD (WENO-Z) and FV (VL2) on the Pallas backend, single precision.
    The per-GPU sub-domain is fixed at $128\times2048\times2048$
    ($5.4\cdot10^{8}$ cells).  Runtimes are the warm execution of 10 (FD)
    or 100 (FV) time steps; efficiency is $T(1)/T(G)$ (ideal $=100\%$).}
    \label{tab:weak_scaling_hydro}
    \begin{tabular}{rrlrrrrrrr}
        \toprule
        & & & & \multicolumn{3}{c}{FD (WENO-Z)} & \multicolumn{3}{c}{FV (VL2)} \\
        \cmidrule(lr){5-7}\cmidrule(lr){8-10}
        GPUs & Nodes & Global grid & Cells & $T$ [s] & Eff.\ [\%] & Mcell/s/GPU & $T$ [s] & Eff.\ [\%] & Mcell/s/GPU \\
        \midrule
        1 & 1 & $128\times2048^{2}$ & $5.4\cdot10^{8}$ & 7.92 & 100.0 & 678 & 7.42 & 100.0 & 7239 \\
        2 & 1 & $256\times2048^{2}$ & $1.1\cdot10^{9}$ & 10.48 & 75.6 & 512 & 11.29 & 65.7 & 4754 \\
        4 & 1 & $512\times2048^{2}$ & $2.1\cdot10^{9}$ & 10.49 & 75.5 & 512 & 11.32 & 65.5 & 4741 \\
        8 & 2 & $1024\times2048^{2}$ & $4.3\cdot10^{9}$ & 11.13 & 71.2 & 482 & 12.05 & 61.6 & 4456 \\
        16 & 4 & $2048^{3}$ & $8.6\cdot10^{9}$ & 11.16 & 71.0 & 481 & 13.78 & 53.8 & 3897 \\
        32 & 8 & $4096\times2048^{2}$ & $1.7\cdot10^{10}$ & 10.92 & 72.5 & 491 & 13.14 & 56.4 & 4086 \\
        64 & 16 & $8192\times2048^{2}$ & $3.4\cdot10^{10}$ & 12.04 & 65.8 & 446 & 14.31 & 51.8 & 3752 \\
        128 & 32 & $16384\times2048^{2}$ & $6.9\cdot10^{10}$ & 11.49 & 68.9 & 467 & 12.82 & 57.9 & 4189 \\
        \bottomrule
    \end{tabular}
\end{table}

\begin{table}
    \centering
    \small
    \setlength{\tabcolsep}{4pt}
    \caption{As Table~\ref{tab:weak_scaling_hydro} but for ideal MHD (the
    CP Alfv\'en-wave setup at benchmark amplitude), per-GPU sub-domain
    $128\times2048\times1024$ ($2.7\cdot10^{8}$ cells).  FD runs on
    several GPUs use the staged Pallas constrained-transport kernels, which
    replace the per-shift collective permutes of the default roll-based
    stencils by one halo exchange per kernel.  The single-GPU rung uses the Pallas
    constrained-transport kernels as well; the default constrained
    transport takes $10.86$~s there.}
    \label{tab:weak_scaling_mhd}
    \begin{tabular}{rrlrrrrrrr}
        \toprule
        & & & & \multicolumn{3}{c}{FD (WENO-Z)} & \multicolumn{3}{c}{FV (VL2)} \\
        \cmidrule(lr){5-7}\cmidrule(lr){8-10}
        GPUs & Nodes & Global grid & Cells & $T$ [s] & Eff.\ [\%] & Mcell/s/GPU & $T$ [s] & Eff.\ [\%] & Mcell/s/GPU \\
        \midrule
        1 & 1 & $128\times2048\times1024$ & $2.7\cdot10^{8}$ & 10.77 & 100.0 & 249 & 10.74 & 100.0 & 2498 \\
        2 & 1 & $256\times2048\times1024$ & $5.4\cdot10^{8}$ & 12.53 & 85.9 & 214 & 14.51 & 74.0 & 1850 \\
        4 & 1 & $512\times2048\times1024$ & $1.1\cdot10^{9}$ & 12.69 & 84.8 & 211 & 14.49 & 74.2 & 1853 \\
        8 & 2 & $1024\times2048\times1024$ & $2.1\cdot10^{9}$ & 13.04 & 82.6 & 206 & 15.31 & 70.2 & 1754 \\
        16 & 4 & $2048\times2048\times1024$ & $4.3\cdot10^{9}$ & 13.17 & 81.8 & 204 & 15.91 & 67.5 & 1687 \\
        32 & 8 & $4096\times2048\times1024$ & $8.6\cdot10^{9}$ & 13.37 & 80.5 & 201 & 15.73 & 68.3 & 1706 \\
        64 & 16 & $8192\times2048\times1024$ & $1.7\cdot10^{10}$ & 13.40 & 80.3 & 200 & 15.81 & 67.9 & 1697 \\
        128 & 32 & $16384\times2048\times1024$ & $3.4\cdot10^{10}$ & 13.19 & 81.6 & 203 & 16.65 & 64.5 & 1612 \\
        \bottomrule
    \end{tabular}
\end{table}

\begin{figure}
    \centering
    \includesvg[width=\linewidth]{figures/weak_scaling.svg}
    \caption{Weak scaling on NVIDIA GH200 GPUs (4 per node) up to 128 GPUs
    on 32 nodes: efficiency $T(1)/T(G)$ (left) and per-GPU throughput
    (right) of the FD (WENO-Z) and FV (VL2) schemes for hydrodynamics and
    MHD (Tables~\ref{tab:weak_scaling_hydro} and~\ref{tab:weak_scaling_mhd}).}
    \label{fig:weak_scaling}
\end{figure}

Within a node all configurations keep $66$--$86\%$ efficiency
(Tables~\ref{tab:weak_scaling_hydro} and~\ref{tab:weak_scaling_mhd},
Fig.~\ref{fig:weak_scaling}).  The onset of the halo exchange costs the
largest step, because a single-GPU step is short (FD hydro: $0.79$~s);
beyond the first node boundary the curves are nearly flat.  At 128 GPUs on 32 nodes FD reaches $68.9\%$ (hydro,
$16384\times2048^{2}$ grid, $6.9\cdot10^{10}$ cells, $467$~Mcell/s/GPU) and
$81.6\%$ (MHD, $203$~Mcell/s/GPU), and FV $57.9\%$ (hydro, $4189$~Mcell/s/GPU) and
$64.5\%$ (MHD, $1612$~Mcell/s/GPU).  The FV scheme updates
$7.2\cdot10^{9}$ (hydro) and $2.5\cdot10^{9}$ (MHD) cells per second on
one GPU, $10.7\times$ and $10.0\times$ the FD throughput; its shorter steps
weigh the halo exchange more, but at 128 GPUs it still sustains
$9.0\times$ the FD throughput for hydrodynamics.

\begin{table}
    \centering
    \caption{Weak scaling for hydrodynamics with self-gravity (the
    Jeans-wave setup at benchmark amplitude, FD with WENO-JS weights, one
    process per GPU), per-GPU sub-domain $64\times1024\times1024$
    ($6.7\cdot10^{7}$ cells), with the slab-decomposed distributed FFT
    Poisson solver and the XLA latency-hiding scheduler.  Numbers in
    brackets give the replicated-FFT baseline, whose per-device memory
    grows with the \emph{global} grid (out of memory beyond 64 GPUs;
    single-execution measurements including first-launch costs).}
    \label{tab:weak_scaling_selfgrav}
    \begin{tabular}{rrlrrrr}
        \toprule
        GPUs & Nodes & Global grid & Cells & Runtime [s] & Eff.\ [\%] & Mcell/s/GPU \\
        \midrule
        1 & 1 & $64\times1024^{2}$ & $6.7\cdot10^{7}$ & 1.69 [1.80] & 100.0 & 397 \\
        2 & 1 & $128\times1024^{2}$ & $1.3\cdot10^{8}$ & 2.03 [2.29] & 83.5 & 331 \\
        4 & 1 & $256\times1024^{2}$ & $2.7\cdot10^{8}$ & 2.04 [3.81] & 82.9 & 329 \\
        8 & 2 & $512\times1024^{2}$ & $5.4\cdot10^{8}$ & 3.16 [4.68] & 53.6 & 213 \\
        16 & 4 & $1024^{3}$ & $1.1\cdot10^{9}$ & 3.53 [7.33] & 47.9 & 190 \\
        32 & 8 & $2048\times1024^{2}$ & $2.1\cdot10^{9}$ & 3.72 [12.33] & 45.5 & 180 \\
        64 & 16 & $4096\times1024^{2}$ & $4.3\cdot10^{9}$ & 4.22 [23.34] & 40.0 & 159 \\
        128 & 32 & $8192\times1024^{2}$ & $8.6\cdot10^{9}$ & 4.58 [---] & 36.9 & 147 \\
        \bottomrule
    \end{tabular}
\end{table}

For self-gravity (Table~\ref{tab:weak_scaling_selfgrav}) the limiting
factor is the FFT Poisson solve.  The XLA SPMD partitioner replicates the
operand of an FFT along its transform dimensions, so the straightforward
formulation costs global-grid memory and work on \emph{every} device --
per-GPU throughput then decays roughly as $1/G$ (bracketed numbers), and
the global grid is capped by single-device memory.  We therefore added a
slab-decomposed distributed Poisson solver (shard-map with all-to-all
transposes between the $x$- and $y$-pencil decompositions, mathematically
identical to the replicated transform), which keeps the per-device
footprint flat ($12.3$~GB at every rung versus $65$~GB at 64~GPUs for the
replicated path) and removes the global-size cap.  Together with the XLA
latency-hiding scheduler (which overlaps the transpose collectives with
stencil compute at no single-GPU cost) the solver reaches $40\%$
efficiency at 64 GPUs -- $5.5\times$ the replicated baseline -- and still
$37\%$ at 128 GPUs on an $8192\times1024^{2}$ grid, sustaining
147~Mcell/s/GPU.  The warm-repeat measurement matters most for this
configuration: constructing the all-to-all communicators costs $0.3$~s on
a single GPU and $4.4$~s at 128 GPUs, so single-execution timings
overstate the ten-step runtime by up to $\sim\!2\times$ at scale.
% ---------------------------------------------------------------------------
